% Generated draft: requires mathematical and source review.
\documentclass[11pt]{article}
\usepackage[margin=25mm]{geometry}
\usepackage{fontspec}
\setmainfont{Latin Modern Roman}
\usepackage{amsmath,amssymb,mathtools,unicode-math}
\setmathfont{Latin Modern Math}
\usepackage{enumitem,needspace,xurl,hyperref}
\setlength{\parindent}{0pt}
\setlength{\parskip}{5pt}
\setlength{\emergencystretch}{3em}
\newcommand{\N}{\mathbb{N}}
\newcommand{\Z}{\mathbb{Z}}
\newcommand{\Q}{\mathbb{Q}}
\newcommand{\R}{\mathbb{R}}
\newcommand{\C}{\mathbb{C}}
\newcommand{\F}{\mathbb{F}}
\newcommand{\A}{\mathbb{A}}
\newcommand{\PP}{\mathbb{P}}
\newcommand{\E}{\mathbb{E}}
\newcommand{\eps}{\varepsilon}
\newcommand{\dd}{\,\mathrm{d}}
\begin{document}
\section*{UnsolvedMath ID: 10400145}
\textbf{Source problem code:} catalogue entry
\subsection{Definitions and mathematical statement}
The notation \(R_1\), its element \(q\), and the invariants \(I_{\mathfrak{sl}_n}\) are inherited from the preceding source discussion, which is absent from this record. Subject to that qualification, let \(t\) be a new indeterminate and put
\[
 J_N=\left(\prod_{j=1}^{N}(t-q^j)\right)\subset R_1[t],\qquad
 R'_1=\varprojlim_{N\geq1}R_1[t]/J_N,
\]
with the transition maps induced by \(J_{N+1}\subseteq J_N\). For each integer \(n\geq1\), evaluation at \(t=q^n\) on any quotient with \(N\geq n\) induces a compatible evaluation map \(\operatorname{ev}_n:R'_1\to R_1\), provided \(R_1\) and \(q\) have the intended coefficient-ring meaning.

\emph{Conditional conjecture (K. Habiro).} Suppose Conjecture 7.29 holds. There exists an invariant \(I_{\mathfrak{sl}}\) of integral homology \(3\)-spheres with values in \(R'_1\) such that, for every such sphere \(M\) and every integer \(n\geq1\),
\[
 \operatorname{ev}_n\bigl(I_{\mathfrak{sl}}(M)\bigr)
 =I_{\mathfrak{sl}_n}(M),\qquad
 I_{\mathfrak{sl}_1}(M)=1.
\]
An integral homology \(3\)-sphere here means a closed connected \(3\)-manifold with integral homology isomorphic to that of \(S^3\); the source excerpt does not specify its orientation convention. The statement of Conjecture 7.29, the precise coefficient ring, the definitions and normalizations of the rank-dependent invariants, and their equivalence convention are not supplied. The conditional hypothesis must therefore remain explicit; no independent existence or uniqueness assertion is substituted for it.
\subsection{Short English statement}
Assuming the preceding Conjecture 7.29, Habiro predicts a single completed-ring-valued invariant of integral homology \(3\)-spheres whose specialization at \(t=q^n\) gives the \(\mathfrak{sl}_n\) invariant for every positive rank, with the rank-one value set to \(1\).
\subsection{Sources}
\begin{itemize}
\item Local UnsolvedMath registry, permanent ID 10400145, source code catalogue entry. Attribution in the supplied statement: K. Habiro; supplied title ``Conjecture 7.30 --- (K.''.
\item Ohtsuki, \emph{Problems on invariants of knots and 3-manifolds} (2002), Conjecture 7.30, PDF page 120, explicitly conditional on Conjecture 7.29. Supplied URL: \url{https://msp.org/gtm/2002/04/gtm-2002-04-024s.pdf}; recorded access 2026-07-29, extraction pdf-font-structure. The collection identifies 231 numbered bold problem, conjecture, and question headings.
\item The supplied registry labels the problem open and OPEN-TRIAGE, describing it as Habiro's unified \(\mathfrak{sl}_n\) invariant and ``Not obviously resolved.'' Source status evidence remains NEEDS\_REVIEW pending a current primary-literature audit. Redistribution terms for the publicly accessible author/publisher copy also require review.
\end{itemize}
\end{document}
